%% expect: native native native native native
\documentclass{article}
\usepackage{tikz-cd}
% A cell that starts with an arrow: \halign's search for \omit expands \arrow, \ar or a shortcut to its \relax on the arrow's line, and
% real PGF's node, with its formula, begins there, whether the cell's text follows on a later line or the cell has none.  LaTeX's
% \check@mathfonts logs the fonts it finds the first time it sets up a size, with that input line.
\makeatletter
\newcommand\state[1]{\typeout{[[#1: \glb@currsize]]}}
\makeatother
\begin{document}
Text.
% The document's first formula, in a cell with no text.
\begin{tikzcd}
\arrow[r] & B
\end{tikzcd}
\state{1}

% Text on the line after the arrow.
\begin{tikzcd}[cells={font=\small}]
\arrow[r]
A & B
\end{tikzcd}
\state{2}

% A label, and the next cell on the next line.
\begin{tikzcd}[cells={font=\large}]
\ar[r, "f"]
  & B
\end{tikzcd}
\state{3}

% The last cell, which \end{tikzcd} ends on the line after it, and a shortcut.
\begin{tikzcd}[cells={font=\Huge}]
|[font=\small]| A \\
\uar
\end{tikzcd}
\state{4}

% With |[...]| before the arrow, the node begins at its closing |.
\begin{tikzcd}
A & |[font=\Large]|
\arrow[l]
\end{tikzcd}
\state{5}
\end{document}
